Extreme value theory (EVT) focuses on the tail of the score distribution. In peaks over threshold (POT), a generalized Pareto distribution is fitted to excesses above an initial fitting threshold; a separate risk parameter specifies the nominal tail probability used to obtain the final alarm threshold \cite{pot_example}. Selecting the fitting threshold involves a bias--variance trade-off: a lower threshold provides more exceedances but may weaken the tail approximation, whereas a higher threshold leaves fewer observations for estimation \cite{scarrott_macdonald_evt_review}. Automatic procedures can select this fitting threshold using goodness-of-fit criteria \cite{automatic_pot_1}; this task is distinct from choosing the target tail probability. SPOT and DSPOT update the fitted tail model as observations arrive, while retaining a user-specified target alarm rate \cite{pot_example}.
